\subsection{分次正则代数}

在本节中，假设 \(H_{0}\) 是一个域，且 H 是有限型的 \(H_{0}\)-代数。

置 \(H_{+} = H_{\geqslant 1}\)；它是 H 的一个极大理想。环 \(S^{-1}H\) 可识别为 H 的局部环 \(H_{H_{+}}\)，对应于理想 \(H_{+}\)；其极大理想是 \((\mathrm{H}_{+})_{\mathrm{H}_{+}} = \mathrm{S}^{-1}\mathrm{H}_{+}\)，其剩余域可识别为 \(H_{0}\)。

命题 8. --- a) 有 \(\dim(\mathbf{H}) \leqslant [\mathbf{H}_{+}/\mathbf{H}_{+}^{2}:\mathbf{H}_{0}]\)。

\begin{enumerate}
\def\labelenumi{\alph{enumi})}
\setcounter{enumi}{1}
\tightlist
\item
  以下条件等价：
\end{enumerate}

\begin{enumerate}
\def\labelenumi{(\roman{enumi})}
\item
  \(\dim (\mathrm{H}) = [\mathrm{H}_{+} / \mathrm{H}_{+}^{2}:\mathrm{H}_{0}]\) ;
\item
  诺特局部环 \(\mathbf{S}^{-1}\mathbf{H}\) 是正则的；
\item
  H 作为 \(H_{0}\)-代数由一族次数 \textgreater{} 0 的、在 \(H_{0}\) 上代数无关的齐次元素生成。
\end{enumerate}

\begin{enumerate}
\def\labelenumi{\alph{enumi})}
\setcounter{enumi}{2}
\tightlist
\item
  假设 b) 的条件成立，并令 \(a_{1}, \ldots, a_{n} \in \mathbf{H}\) 是次数 \textgreater{} 0 的齐次元素。以下条件等价：
\end{enumerate}

\begin{enumerate}
\def\labelenumi{(\roman{enumi})}
\item
  \(a_{i}\) 的类构成 \(\mathbf{H}_{0}\)-向量空间 \(\mathbf{H}_{+}/\mathbf{H}_{+}^{2}\) 的一个基；
\item
  \(a_{i}/1\) 构成正则诺特局部环 \(\mathbf{S}^{-1}\mathbf{H}\) 的一个坐标系；
\item
  \(a_{i}\) 在 \(\mathbf{H}_{0}\) 上代数无关，并生成 \(\mathbf{H}\)（作为 \(\mathbf{H}_{0}\)-代数）。
\end{enumerate}

我们有 \(\dim(\mathbf{H})=\dim(\mathbf{S}^{-1}\mathbf{H})\)（第 2 号，定理 1），并且 \([\mathbf{H}_{+}/\mathbf{H}_{+}^{2}:\mathbf{H}_{0}]=[(\mathbf{S}^{-1}\mathbf{H}_{+})/(\mathbf{S}^{-1}\mathbf{H}_{+})^{2}:\mathbf{H}_{0}]\)（II, § 3, 第 3 号，命题 9）；根据 § 5，第 1 号，这推出 \(a\) 以及等价关系 (i) \(\Leftrightarrow\) (ii)（分别在 \(b\) 和 \(c\) 中）。其中的两个蕴含 (iii) \(\Rightarrow\) (i) 分别位于 \(b\) 和 \(c\) 中，显然成立。现证明蕴含 (i) \(\Rightarrow\) (iii)。因此，假设 \(\dim(\mathbf{H})=[\mathbf{H}_{+}/\mathbf{H}_{+}^{2}:\mathbf{H}_{0}]\)，并令齐次元素 \(a_{1},\ldots,a_{n}\) 属于 H，次数为 \(>0\)，且它们的类构成 \(\mathbf{H}_{+}/\mathbf{H}_{+}^{2}\) 的一个基；考虑分次代数同态 \(\varphi:\mathrm{H}_{0}[\mathrm{X}_{1},\ldots,\mathrm{X}_{n}]\to\mathrm{H}\)，它将 \(\mathrm{X}_{i}\) 送到 \(a_{i}\)。H 的理想 \(\mathrm{H}_{+}\) 由 \(a_{i}\) 生成（A, II, p.~171, 推论 2 和注）；因此同态 \(\varphi\) 是满射（III, § 1, 第 2 号，命题 1）。但是，有

\[
\dim \left(\mathrm{H} _ {0} \left[ \mathrm{X} _ {1}, \dots , \mathrm{X} _ {n} \right]\right) = n = \dim (\mathrm{H}) = \dim \left(\mathrm{H} _ {0} \left[ \mathrm{X} _ {1}, \dots , \mathrm{X} _ {n} \right] / (\operatorname{Ker} \varphi)\right)
\]

（§ 2，第 4 号，定理 3 的推论 1），因此 Ker φ = 0，因为 H₀{[}X₁, \ldots, Xₙ{]} 是整环。这就证明了断言 (iii)。

在 \(b\) ) 的假设下，称 H 为正则的 \(\mathrm{H}_{0}\)-分次代数，或 \(\mathrm{H}_{0}\)-分次多项式代数。在 \(c\) ) 的假设下，称 \(a_{i}\) 构成 H 的一个分次坐标系。

注记。--- 1) 采用 \(c\) ) 的记号，令 \(d_{i}\) 为 \(a_{i}\) 的次数（ \(1 \leqslant i \leqslant n\) ）。则庞加莱级数 \(\mathbf{P}_{\mathrm{H}} = \sum_{n \in \mathbf{Z}} [\mathrm{H}_{n} : \mathrm{H}_{0}] \cdot \mathrm{T}^{n}\) 等于 \(\prod_{i} (1 - \mathrm{T}^{d_{i}})^{-1}\)（§4，第 2 号，例 1）。因此，若 H 是 \(\mathrm{H}_{0}\)-分次多项式代数，则有

\[
\mathrm{P} _ {\mathrm{H}} = \prod_ {n} (1 - \mathrm{T} ^ {n}) ^ {- \delta (n)}, \quad \text { with } \quad \delta (n) = \left[ (\mathrm{H} _ {+} / \mathrm{H} _ {+} ^ {2}) _ {n}: \mathrm{H} _ {0} \right].
\]

\begin{enumerate}
\def\labelenumi{\arabic{enumi})}
\setcounter{enumi}{1}
\tightlist
\item
  反之，存在整数 \(d_{i}>0\) 使 \(\mathbf{P}_{\mathbf{H}}=\prod_{i}(1-\mathbf{T}^{d_{i}})^{-1}\)，并不意味着 H 是分次多项式代数；例如，若 H 由一个次数为 1 的元素 X 和一个次数为 2 的元素 Y 生成，且唯一关系为 \(\mathbf{X}^{2}=0\)，则有 \(\mathbf{P}_{\mathbf{H}}=(1-\mathbf{T})^{-1}\)。
\end{enumerate}

